% Preamble
\documentclass{article}

% Packages
\usepackage{hyperref}

\title{Idei De bază pentru teza de Dizertație}
\author{Marius-Vlad Dumitru}
% Document
\begin{document}
\maketitle
\section{Motivație. Punctul de plecare inițial.}
Am pornit de la ideea de Reinforcement Learning/Deep Reinforcement Learning.

La BIOSINF am studiat mai putin despre RL/DRL, iar pentru ca educația primita de la masterul BIOSINF să fie completă("Round Trip") mi s-a părut interesant să aleg o temă de disertație bazată de RL/DRL.

Am fost mereu pasionat de jocurile video, iar RL/DRL se pretează pentru așa ceva, deci teza de disertație este oportunitatea care permite studiul RL/DRL aplicat pentru jocuri video.

Prin urmare am intenționat să dezvolt agenți RL/DRL pentru un joc video.

Prima provocare a fost alegera unui joc.

\section{Alegerea Jocului}
M-am orientat către un joc de strategie, de tip RTS(Real Time Strategy), întrucât acest tip de joc se pretează pentru studiul RL/DRL, iar personal am jucat câteva jocuri de strategie precum Starcraft, Starcraft 2, Warcraft 3, Age of Empires, etc.

Inițial am vrut să folosesc unul din jocurile de la Blizzard, fie Starcraft, Starcraft 2 sau Warcraft 3. Din păcate niciunul din ele nu oferă o interfață nativă în Python pentru controlul extern al jucătorilor.

Starcraft 2 oferă o interfață pentru controlul extern al jucătorilor cu ajutorul proiectului [S2Client-Proto](https://github.com/blizzard/s2client-proto).

Starcraft Broodwar oferă o interfață pentru controlul extern al jucătorilor cu ajutorul proiectului [BWAPI](https://bwapi.github.io/) dar acesta funcționează numai pentru versiunea 1.6.

Deși mi-ar fi plăcut foarte mult să folosesc Starcraft/Starcraft 2/Warcraft 3 pentru realizarea experimentelor, am realizat că aceste jocuri sunt mult prea complexe pentru nivelul meu de cunoștințe în RL/DRL. Prin urmare am fost nevoit să mă orientez către altceva.

Jocul pe care l-am ales să îl folosesc în experimente este MicroRTS, versiunea originală disponibilă \href{https://github.com/Farama-Foundation/MicroRTS}{AICI}.

Am ales acest joc din următoarele motive:
\begin{enumerate}
\item \textbf{Este mai simplu din punct de vedere al RL/DRL. Nu este la fel de complex precum Starcraft/Starcraft 2/Warcraft 3}. Fiind mai simplu, este mai ușor de lucrat cu el și posibilitatea de a greși este mai mică. Spațiile alea (acțiuni, stări, etc.) sunt mult mai mici față de Warcraft 3/Starcraft/Starcraft2.
\item Deși mai simplu decât jocurile Blizzard, este "good enough" de complex pentru experimente de tip RL/DRL.
\item Există paper-uri publicate folosing MicroRTS.
\item Este open-source, deci codul sursă este disponibil și este permisă modificarea și redistribuția sa.
\item Are deja boți predefiniți în joc pe care îi pot folosi.
\item Este special construit pentru experimente de tip RL/DRL.
\end{enumerate}


% Some ideas that i don't want to forget
%Aici e ideile mari boss
%uite niste idei noi:
%Suna misto:
%Win Condition/ Win conditions
%Constraints(constrangeri)
\end{document}